% ============================================================
%  参考文献、译名对照、原文勘误
% ============================================================
\raggedright

\section*{附录 A：参考文献}
\addcontentsline{toc}{section}{附录 A：参考文献}

\subsection*{A 级：原始文献、经典教材、作者本人文献}
\begin{enumerate}
\item S.~Raskin (lecturer), L.~Chen (notes), \emph{Notes for ``Geometric Langlands for Projective Curves''},
July 07--11, 2025, 18 pp. 公开副本：
\url{https://www.math.columbia.edu/~okounkov/Sam_lecture_notes.pdf}.
\item V.~Drinfeld, D.~Gaitsgory, \emph{Geometric constant term functor(s)},
Selecta Math. (N.S.) \textbf{22} (2016), 1881--1951. Zbl 1360.14060.
\item B.~C.~Ngo, P.~Polo, \emph{Résolutions de Demazure affines et formule de Casselman--Shalika géométrique},
arXiv:math/0005022（预印本，承诺发表于 J. Algebraic Geom.；底本 §5.1 将公式归给 Frenkel--Gaitsgory--Vilonen）.
\item E.~Frenkel, D.~Gaitsgory, K.~Vilonen, \emph{Whittaker patterns in the geometry of moduli spaces of bundles on curves},
Ann. of Math. (2) \textbf{153} (2001), 699--748.
\item E.~Frenkel, D.~Gaitsgory, K.~Vilonen, \emph{On the geometric Langlands conjecture},
J. Amer. Math. Soc. \textbf{15} (2002), 367--417; arXiv:math/0012255.
\item D.~Gaitsgory, S.~Raskin, \emph{Proof of the geometric Langlands conjecture I: construction of the functor},
arXiv:2405.03599.
\item D.~Arinkin, D.~Gaitsgory, \emph{Singular support of coherent sheaves and the geometric Langlands conjecture},
Selecta Math. (N.S.) \textbf{21} (2015), 1--199.
\item P.~Deligne, \emph{Équations différentielles à points singuliers réguliers},
Lecture Notes in Math. \textbf{163}, Springer, 1970（Riemann--Hilbert 对应与正则奇点）.
\end{enumerate}

\subsection*{B 级：可核对的百科、正式讲义}
\begin{enumerate}
\item T.~Feng, \emph{Geometric class field theory}, 讲义,
\url{https://math.berkeley.edu/~fengt/2GeometricCFT.pdf}.
\item nLab, ``local system'', Last revised 2025-04-10,
\url{https://ncatlab.org/nlab/show/local+system}.
\item nLab, ``perverse sheaf'', \url{https://ncatlab.org/nlab/show/perverse+sheaf}.
\item nLab, ``geometric Langlands correspondence'',
\url{https://ncatlab.org/nlab/show/geometric+Langlands+correspondence}.
\item nLab, ``Ran space'', \url{https://ncatlab.org/nlab/show/Ran+space}.
\item nLab, ``Whittaker function'', \url{https://ncatlab.org/nlab/show/Whittaker+function}.
\end{enumerate}

\subsection*{I 级：本库既有译注、术语表、资料卡}
\begin{enumerate}
\item 本库既有卡 \texttt{glc.concept.deligne-gl1-existence}（Deligne GL\(_1\) 定理，本质唯一）.
\item 本库既有卡 \texttt{glc.concept.unramified-global-gl}（非分歧整体几何 Langlands 猜想）.
\item 本库既有卡 \texttt{glc.concept.cuspidal-de-rham}（尖点 de Rham 等价）.
\item 本库既有卡 \texttt{gaits04.concept.eigensheaf}（Hecke 特征层）.
\item 本库既有卡 \texttt{bdh.concept.hecke-eigensheaf}（D-模侧的 Hecke 特征层）.
\end{enumerate}

\section*{附录 B：译名对照}
\addcontentsline{toc}{section}{附录 B：译名对照}

\begin{center}
\small
\begin{tabular}{lll}
\hline
英文 & 中文 & 备注 \\
\hline
local system & 局部系统 & 与 lisse 层同义 \\
lisse & lisse & 保留原文，不译 \\
perverse sheaf & 反常层 & 严禁「错致层」 \\
cuspidal & 尖点 & 作形容词时写「尖点层」 \\
constant term & 常数项 & \\
Whittaker coefficient & Whittaker 系数 & 人名不译 \\
Hecke eigensheaf & Hecke 特征层 & \\
geometric Satake equivalence & 几何 Satake 等价 & \\
affine Grassmannian & 仿射 Grassmann 簇 & \\
loop group & 圈群 & \\
positive loop group & 正圈群 & \\
local Hecke stack & 局部 Hecke 叠 & \\
ind-coherent sheaves & 归纳凝聚层 & IndCoh 保留记号 \\
perfect complex & 完美复形 & \\
quasicoherent & 拟凝聚 & \\
six-functor formalism & 六函子形式化 & \\
correspondence & 对应 & \\
constructible sheaf & 可构造层 & \\
flat connection & 平坦联络 & \\
D-module & D-模 & \\
nilpotent singular support & 幂零奇异支撑 & \\
singular support & 奇异支撑 & \\
microlocal & 微局部 & \\
microstalk & 微茎 & \\
cotangent bundle & 余切丛 & \\
Borel subgroup & Borel 子群 & 人名不译 \\
maximal torus & 极大环面 & \\
reductive group & 约化群 & \\
unipotent radical & 单幂根 & \\
parabolic subgroup & 抛物子群 & \\
Levi subgroup & Levi 子群 & \\
Langlands dual group & Langlands 对偶群 & \\
Eisenstein series & Eisenstein 级数 & \\
Kac--Moody localization & Kac--Moody 局部化 & \\
critical level & 临界水平 & \\
spherical vector & 球向量 & \\
automorphic form & 自守形式 & \\
automorphic representation & 自守表示 & \\
function field & 函数域 & \\
\hline
\end{tabular}
\end{center}

\section*{附录 C：原文勘误表}
\addcontentsline{toc}{section}{附录 C：原文勘误表}

下表只列本译本核对过的、影响理解的原文笔误；拼写类错误集中在末行，不逐条列入。

\begin{center}
\small
\begin{tabular}{>{\raggedright\arraybackslash}p{1.2cm}>{\raggedright\arraybackslash}p{4.4cm}>{\raggedright\arraybackslash}p{3.8cm}>{\raggedright\arraybackslash}p{3.8cm}}
\hline
页 & 原文 & 应为 & 依据 \\
\hline
4 & \(AJ_d\) 的纤维是 \(\Gamma(X,L)\setminus 0\) & \(\mathbb P\Gamma(X,L)\) & 按除子定义；Feng 讲义 \\
4 & preverse & perverse & 拼写 \\
6 & \(B H^1(X,L)\)；\(H^1(X,L)\) 可缩 & \(R\Gamma(X,L)[1]\)；\(d>2g-2\) 时为 \(B H^0(X,L)\) & 脚注与 \(H^1=0,H^0\neq0\) \\
6 & \(0\to L\to B\to \cO\to 0\) & \(0\to L\to \cE\to \cO\to 0\) & 正合列中 \(B\) 为 \(\cE\) 之误 \\
7 & holonomic modular form & holomorphic modular form & 上下文 \\
9 & \(\sigma^{(d)}\) perverse but not lisse & \(\sigma^{(d)}\) lisse；\(\cF_\sigma\) perverse but not lisse & 定义 \\
12 & \(\Rep(G)\to\Shv(X)\) & \(\Rep(\Gd)\to\Shv(X)\) & \(\Gd\)-局部系统的定义 \\
16 & Vilonon & Vilonen & 作者名 \\
多处 & irrducible, uniquelly, Geothendieck, numbder, restiction, Frobenious, archemedean, mearsure, swaping, wher, cmoplexes, prametrizing & 对应正字 & 拼写 \\
\hline
\end{tabular}
\end{center}

\section*{附录 D：底本与版权说明}
\addcontentsline{toc}{section}{附录 D：底本与版权说明}

底本为 Sam Raskin 讲授、Lin Chen 笔记的公开讲义，原 PDF 共 18 页，PDF 页与印刷页一致。
本译本是中文学术译注，不替代原文；原文著作权归作者所有。译本中的译注、术语说明与勘误表
由译者撰写，按本库规范以「改写」（rewritten）方式登记。引用本译本时请同时引用底本。
